% GENERATED FILE. Edit canonical lessons, then run npm run generate:books.
% canonical-source-hash: fnv1a64:bde6137673d1c010
% canonical-lessons: GU-C415-cokdi, GU-C415-khaco, GU-C415-nisani, GU-C415-tavar, GU-C415-gotvu

\chapter{Landmarks}
\label{ch:gu-landmarks}
\hlchaptermodality{415}

\begin{chapteropening}
\textbf{By the end of this chapter:} \emph{I can give directions by landmarks.}

Complete the last lesson of chapter 415: I can give directions by landmarks.
\end{chapteropening}

\section[\texorpdfstring{\gu{ચોકડી} (cokḍī)}{ચોકડી (cokḍī)}]{\gu{ચોકડી} (cokḍī) — a crossroads (on a highway)}

\label{lesson:GU-C415-cokdi}



\begin{quote}
Before the new one: say the Gujarati for the wrong way round, then the Gujarati for a compass.
\end{quote}

\subsection*{You'll want to know: \gu{ચોકડી}}
\textbf{\gu{ચોકડી}} — \emph{cokḍī} — \textquotedblleft{}a crossroads (on a highway)\textquotedblright{}.

\begin{quote}
\gu{ચોકડી} \gu{પરથી} \gu{જમણી} \gu{બાજુ} \gu{વળો}.
\end{quote}

\begin{quote}
\emph{cokḍī parthī jamṇī bāju vaḷo.}
\end{quote}

\begin{quote}
Turn right at the crossroads.
\end{quote}

A feminine noun.

\begin{cousinweb}
\textbf{\gu{ચોકડી}} (\emph{cokḍī}) is built on \gu{ચોક}, a square, from Sanskrit \emph{catuṣka}, \textquotedblleft{}a set of four\textquotedblright{} (CDIAL 4629): the place where four roads meet.
\end{cousinweb}

\begin{grammarlens}[title={What you've built}]
One more word for getting around.
\end{grammarlens}

\subsection*{Your turn}
\begin{itemize}
\raggedright
  \item \emph{Say it:} \emph{cokḍī}
  \item \emph{Say it:} \emph{cokḍī}, once more
  \item \emph{Say it:} \emph{cokḍī parthī jamṇī bāju vaḷo.} — Turn right at the crossroads.
  \item \emph{Recall:} say the Gujarati for the wrong way round, then the Gujarati for a compass, then say \emph{cokḍī} again
\end{itemize}

\begin{tcolorbox}[breakable,colback=teal!4,colframe=teal!35!black,title={Before you move on}]
What does \gu{ચોકડી} mean? (\textquotedblleft{}A crossroads (on a highway)\textquotedblright{}.) Say it once more.
\end{tcolorbox}
\section[\texorpdfstring{\gu{ખાંચો} (khā̃co)}{ખાંચો (khā̃co)}]{\gu{ખાંચો} (khā̃co) — a side lane; a nook}

\label{lesson:GU-C415-khaco}



\begin{quote}
Before the new one: say the Gujarati for a compass, then the Gujarati for a crossroads.
\end{quote}

\subsection*{You'll want to know: \gu{ખાંચો}}
\textbf{\gu{ખાંચો}} — \emph{khā̃co} — \textquotedblleft{}a side lane; a nook\textquotedblright{}.

\begin{quote}
\gu{દુકાન} \gu{પેલા} \gu{ખાંચામાં} \gu{છે}.
\end{quote}

\begin{quote}
\emph{dukān pelā khā̃cāmā̃ che.}
\end{quote}

\begin{quote}
The shop is in that side lane.
\end{quote}

A masculine noun.

\begin{cousinweb}
\textbf{\gu{ખાંચો}} (\emph{khā̃co}) may belong with a Gujarati word for a corner, which CDIAL lists under \emph{*khōñcā}, \textquotedblleft{}a curve, a bend\textquotedblright{}, a form it reconstructs (CDIAL 3930); CDIAL gives no form for \gu{ખાંચો} itself, so the link is hedged.
\end{cousinweb}

\begin{grammarlens}[title={What you've built}]
One more word for getting around.
\end{grammarlens}

\subsection*{Your turn}
\begin{itemize}
\raggedright
  \item \emph{Say it:} \emph{khā̃co}
  \item \emph{Say it:} \emph{khā̃co}, once more
  \item \emph{Say it:} \emph{dukān pelā khā̃cāmā̃ che.} — The shop is in that side lane.
  \item \emph{Recall:} say the Gujarati for a compass, then the Gujarati for a crossroads, then say \emph{khā̃co} again
\end{itemize}

\begin{tcolorbox}[breakable,colback=teal!4,colframe=teal!35!black,title={Before you move on}]
What does \gu{ખાંચો} mean? (\textquotedblleft{}A side lane; a nook\textquotedblright{}.) Say it once more.
\end{tcolorbox}
\section[\texorpdfstring{\gu{નિશાની} (niśānī)}{નિશાની (niśānī)}]{\gu{નિશાની} (niśānī) — a landmark; a sign}

\label{lesson:GU-C415-nisani}



\begin{quote}
Before the new one: say the Gujarati for a crossroads, then the Gujarati for a side lane.
\end{quote}

\subsection*{You'll want to know: \gu{નિશાની}}
\textbf{\gu{નિશાની}} — \emph{niśānī} — \textquotedblleft{}a landmark; a sign\textquotedblright{}.

\begin{quote}
\gu{અમારા} \gu{ઘરની} \gu{નિશાની} \gu{એક} \gu{મોટો} \gu{લીમડો} \gu{છે}.
\end{quote}

\begin{quote}
\emph{amārā gharnī niśānī ek moṭo līmḍo che.}
\end{quote}

\begin{quote}
The landmark for our house is a big neem tree.
\end{quote}

A feminine noun.

\begin{cousinweb}
\textbf{\gu{નિશાની}} (\emph{niśānī}) is usually traced to Persian \emph{nišānī}, \textquotedblleft{}a mark, a token\textquotedblright{}, built on \emph{nišān}, \textquotedblleft{}a sign\textquotedblright{}: something you know a place by.
\end{cousinweb}

\begin{grammarlens}[title={What you've built}]
One more word for getting around.
\end{grammarlens}

\subsection*{Your turn}
\begin{itemize}
\raggedright
  \item \emph{Say it:} \emph{niśānī}
  \item \emph{Say it:} \emph{niśānī}, once more
  \item \emph{Say it:} \emph{amārā gharnī niśānī ek moṭo līmḍo che.} — The landmark for our house is a big neem tree.
  \item \emph{Recall:} say the Gujarati for a crossroads, then the Gujarati for a side lane, then say \emph{niśānī} again
\end{itemize}

\begin{tcolorbox}[breakable,colback=teal!4,colframe=teal!35!black,title={Before you move on}]
What does \gu{નિશાની} mean? (\textquotedblleft{}A landmark; a sign\textquotedblright{}.) Say it once more.
\end{tcolorbox}
\section[\texorpdfstring{\gu{ટાવર} (ṭāvar)}{ટાવર (ṭāvar)}]{\gu{ટાવર} (ṭāvar) — a clock tower; a tower}

\label{lesson:GU-C415-tavar}



\begin{quote}
Before the new one: say the Gujarati for a side lane, then the Gujarati for a landmark.
\end{quote}

\subsection*{You'll want to know: \gu{ટાવર}}
\textbf{\gu{ટાવર}} — \emph{ṭāvar} — \textquotedblleft{}a clock tower; a tower\textquotedblright{}.

\begin{quote}
\gu{શહેરના} \gu{ટાવરની} \gu{ઘડિયાળ} \gu{બંધ} \gu{છે}.
\end{quote}

\begin{quote}
\emph{śahernā ṭāvarnī ghaḍiyāḷ bandh che.}
\end{quote}

\begin{quote}
The clock on the town's clock tower has stopped.
\end{quote}

A masculine noun.

\begin{cousinweb}
\textbf{\gu{ટાવર}} (\emph{ṭāvar}) is borrowed from English \emph{tower}; in many Gujarati towns the \gu{ટાવર} is the clock tower in the main square.
\end{cousinweb}

\begin{grammarlens}[title={What you've built}]
One more word for getting around.
\end{grammarlens}

\subsection*{Your turn}
\begin{itemize}
\raggedright
  \item \emph{Say it:} \emph{ṭāvar}
  \item \emph{Say it:} \emph{ṭāvar}, once more
  \item \emph{Say it:} \emph{śahernā ṭāvarnī ghaḍiyāḷ bandh che.} — The clock on the town's clock tower has stopped.
  \item \emph{Recall:} say the Gujarati for a side lane, then the Gujarati for a landmark, then say \emph{ṭāvar} again
\end{itemize}

\begin{tcolorbox}[breakable,colback=teal!4,colframe=teal!35!black,title={Before you move on}]
What does \gu{ટાવર} mean? (\textquotedblleft{}A clock tower; a tower\textquotedblright{}.) Say it once more.
\end{tcolorbox}
\section[\texorpdfstring{\gu{ગોતવું} (gotvũ)}{ગોતવું (gotvũ)}]{\gu{ગોતવું} (gotvũ) — to hunt for; to search (everyday)}

\label{lesson:GU-C415-gotvu}



\begin{quote}
Before the new one: say the Gujarati for a landmark, then the Gujarati for a clock tower.
\end{quote}

\subsection*{You'll want to know: \gu{ગોતવું}}
\textbf{\gu{ગોતવું}} — \emph{gotvũ} — \textquotedblleft{}to hunt for; to search (everyday)\textquotedblright{}.

\begin{quote}
\gu{અમે} \gu{અડધો} \gu{કલાક} \gu{સરનામું} \gu{ગોત્યું}.
\end{quote}

\begin{quote}
\emph{ame aḍdho kalāk sarnāmũ gotyũ.}
\end{quote}

\begin{quote}
We hunted for the address for half an hour.
\end{quote}

Past, with \gu{તેણે} (he or she …): \gu{તેણે} \gu{ઘર} \gu{ગોત્યું}, \emph{teṇe ghar gotyũ}.

\begin{cousinweb}
\textbf{\gu{ગોતવું}} (\emph{gotvũ}) No dictionary we use records an older form, so its history is not settled.
\end{cousinweb}

\begin{grammarlens}[title={What you've built}]
One more word for getting around.
\end{grammarlens}

\subsection*{Your turn}
\begin{itemize}
\raggedright
  \item \emph{Say it:} \emph{gotvũ}
  \item \emph{Say it:} \emph{gotvũ}, once more
  \item \emph{Say it:} \emph{ame aḍdho kalāk sarnāmũ gotyũ.} — We hunted for the address for half an hour.
  \item \emph{Recall:} say the Gujarati for a landmark, then the Gujarati for a clock tower, then say \emph{gotvũ} again
\end{itemize}

\begin{tcolorbox}[breakable,colback=teal!4,colframe=teal!35!black,title={Before you move on}]
What does \gu{ગોતવું} mean? (\textquotedblleft{}To hunt for; to search (everyday)\textquotedblright{}.) Say it once more.
\end{tcolorbox}
